\documentclass[12pt, letterpaper]{article}

\usepackage[utf8]{inputenc}
\usepackage[T1]{fontenc}
\usepackage[spanish]{babel}
\usepackage{geometry}
\usepackage{amsmath, amssymb, amsthm}
\usepackage{booktabs}
\usepackage{array}
\usepackage{microtype}
\usepackage{setspace}
\usepackage{parskip}
\usepackage{titlesec}
\usepackage{fancyhdr}
\usepackage{enumitem}
\usepackage{bm}
\usepackage{mathtools}
\usepackage{tcolorbox}
\usepackage{hyperref}

\tcbuselibrary{skins, breakable}

\geometry{top=2.8cm, bottom=3.0cm, left=3.2cm, right=3.2cm}
\setstretch{1.20}

\pagestyle{fancy}
\fancyhf{}
\fancyhead[C]{\small\textsc{Econometria · Gujarati --- Ejercicio 11.4}}
\fancyfoot[C]{\small\thepage}
\renewcommand{\headrulewidth}{0.4pt}
\renewcommand{\footrulewidth}{0pt}

\titleformat{\section}
  {\large\bfseries\scshape}{\thesection.}{0.6em}{}
  [\vspace{-4pt}\rule{\linewidth}{0.4pt}]
\titleformat{\subsection}{\normalsize\bfseries}{\thesubsection.}{0.5em}{}
\titleformat{\subsubsection}{\normalsize\itshape}{}{0em}{}

\newtheorem{prop}{Proposicion}
\newtheorem{lema}{Lema}
\theoremstyle{remark}
\newtheorem{obs}{Observacion}

\newtcolorbox{clave}{
  colback=white, colframe=black,
  boxrule=0.5pt, arc=3pt,
  left=8pt, right=8pt, top=6pt, bottom=6pt,
  breakable
}

\newcommand{\E}{\mathbb{E}}
\newcommand{\Var}{\operatorname{Var}}

\begin{document}

\begin{center}
  {\LARGE\bfseries Modelos Logaritmicos, Distribucion Log-Normal}\\[6pt]
  {\LARGE\bfseries y Propiedades del Termino de Perturbacion}\\[12pt]
  {\large\itshape Distribucion de $u_i$, desigualdad de Jensen y correccion del intercepto}\\[14pt]
  \rule{0.55\linewidth}{0.6pt}\\[8pt]
  {\normalsize Ejercicio~11.4 --- \textit{Econometria}, Gujarati \& Porter, 5.a ed.}\\[4pt]
  {\small\textsc{Con derivacion de la funcion generadora de momentos log-normal,
  demostracion de Jensen y reparametrizacion del intercepto}}
\end{center}

\vspace{10pt}

\noindent\textbf{Planteamiento.}\quad
Considere el modelo de potencia:
\begin{equation}
  Y_i = \beta_1 X_i^{\beta_2} u_i
  \tag{1}\label{eq:nivel}
\end{equation}
que al tomar logaritmos naturales se convierte en:
\begin{equation}
  \ln Y_i = \ln\beta_1 + \beta_2 \ln X_i + \ln u_i
  \tag{2}\label{eq:log}
\end{equation}

Definiendo $\alpha_1 = \ln\beta_1$ y $\varepsilon_i = \ln u_i$, el modelo es:
\begin{equation}
  \ln Y_i = \alpha_1 + \beta_2 \ln X_i + \varepsilon_i
  \label{eq:log_v2}
\end{equation}

Para que MCO sea aplicable a~( \ref{eq:log_v2} ), se necesita $\E[\varepsilon_i] = \E[\ln u_i] = 0$.
Se analizan las implicaciones de este supuesto.

\section{Inciso (a): distribucion de $u_i$ si $\E(\ln u_i) = 0$}

\subsection{Planteamiento}

Se supone que $\varepsilon_i = \ln u_i$ tiene media cero:
\begin{equation}
  \E[\ln u_i] = 0
  \label{eq:cond_a}
\end{equation}

¿Que distribucion debe seguir $u_i$?

\subsection{Distribucion log-normal}

Si se impone ademas que $\varepsilon_i = \ln u_i$ sigue una distribucion normal:
\begin{equation}
  \ln u_i \sim \mathcal{N}(0, \sigma^2)
  \label{eq:lognormal_cond}
\end{equation}
entonces $u_i$ sigue una \textbf{distribucion log-normal} con parametros
$\mu = 0$ y $\sigma^2$. La funcion de densidad de probabilidad de $u_i$ es:
\begin{equation}
  f_{u_i}(x) = \frac{1}{x\sigma\sqrt{2\pi}} \exp\!\left(-\frac{(\ln x)^2}{2\sigma^2}\right),
  \qquad x > 0
  \label{eq:fdp_lognormal}
\end{equation}

\subsection{Momentos de la distribucion log-normal}

Los momentos de $u_i$ se derivan de la funcion generadora de momentos (FGM)
de la distribucion normal. Si $Z = \ln u_i \sim \mathcal{N}(\mu, \sigma^2)$, entonces
$u_i = e^Z$ y:

\begin{equation}
  \E[u_i^r] = \E[e^{rZ}] = e^{r\mu + r^2\sigma^2/2}
  \label{eq:momentos_lognormal}
\end{equation}

Para el primer momento ($r = 1$) con $\mu = 0$:
\begin{equation}
  \E[u_i] = e^{0 + \sigma^2/2} = e^{\sigma^2/2} > 1
  \label{eq:media_u}
\end{equation}

Para el segundo momento ($r = 2$) con $\mu = 0$:
\begin{equation}
  \E[u_i^2] = e^{2\sigma^2}
  \label{eq:mom2_u}
\end{equation}

La varianza es:
\begin{equation}
  \Var(u_i) = \E[u_i^2] - [\E[u_i]]^2
            = e^{2\sigma^2} - e^{\sigma^2}
            = e^{\sigma^2}(e^{\sigma^2} - 1)
  \label{eq:var_u}
\end{equation}

\subsection{Consecuencias}

Bajo~(\ref{eq:lognormal_cond}), aunque $\E[\ln u_i] = 0$, la media del error
en niveles \emph{no es uno sino}:
\begin{equation}
  \E[u_i] = e^{\sigma^2/2} \neq 1 \quad \text{(a menos que } \sigma^2 = 0\text{)}
  \label{eq:media_u_no_uno}
\end{equation}

Esto implica que el modelo en niveles~(\ref{eq:nivel}) \textbf{no cumple}
$\E[u_i] = 1$ (condicion que garantizaria que el error tenga media cero en
la escala multiplicativa). La transformacion logaritmica y la condicion de media
cero en logaritmos son compatibles con la log-normalidad, pero implican que
la media del error en niveles es $> 1$.

\begin{clave}
  Si $\E[\ln u_i] = 0$ y se asume normalidad de $\ln u_i$, entonces $u_i$
  sigue una distribucion \textbf{log-normal} con parametros $(0, \sigma^2)$.
  Su media es $\E[u_i] = e^{\sigma^2/2} > 1$ y su varianza es
  $e^{\sigma^2}(e^{\sigma^2}-1)$.
\end{clave}

\section{Inciso (b): ¿es $\E(\ln u_i) = 0$ si $\E(u_i) = 1$?}

\subsection{Respuesta}

\textbf{No}. Si $\E(u_i) = 1$, entonces $\E(\ln u_i) < 0$.

\subsection{Demostracion mediante la desigualdad de Jensen}

La funcion logaritmo $g(x) = \ln x$ es \textbf{estrictamente cóncava}
para $x > 0$: su segunda derivada es $g''(x) = -1/x^2 < 0$.

\begin{lema}[Desigualdad de Jensen para funciones concavas]
  Si $g$ es estrictamente concava y $X$ es una variable aleatoria no degenerada
  (con varianza positiva), entonces:
  \begin{equation}
    \E[g(X)] < g(\E[X])
    \label{eq:Jensen}
  \end{equation}
\end{lema}

Aplicando~(\ref{eq:Jensen}) con $g(x) = \ln x$ y $X = u_i$, bajo el supuesto
de que $u_i$ no es constante ($\Var(u_i) > 0$):
\begin{equation}
  \E[\ln u_i] < \ln \E[u_i]
  \label{eq:Jensen_aplicado}
\end{equation}

Si $\E[u_i] = 1$:
\begin{equation}
  \E[\ln u_i] < \ln 1 = 0
  \label{eq:Jensen_resultado}
\end{equation}

Por lo tanto, $\E[\ln u_i] < 0$: la media del logaritmo del error es
\textbf{estrictamente negativa} cuando la media del error en niveles es~1.

\subsection{Cuantificacion del sesgo bajo log-normalidad}

Si $u_i \sim \text{Log-Normal}(\mu, \sigma^2)$ (es decir, $\ln u_i \sim \mathcal{N}(\mu, \sigma^2)$),
la condicion $\E[u_i] = 1$ implica:
\begin{equation}
  e^{\mu + \sigma^2/2} = 1 \quad \Rightarrow \quad \mu = -\frac{\sigma^2}{2}
  \label{eq:mu_condicion}
\end{equation}

Por tanto, bajo la condicion $\E[u_i] = 1$:
\begin{equation}
  \E[\ln u_i] = \mu = -\frac{\sigma^2}{2} < 0
  \label{eq:media_lnu_negativa}
\end{equation}

La magnitud del sesgo ($\sigma^2/2$) aumenta con la varianza del logaritmo
del error. Para varianzas pequenas el sesgo es pequeno, pero para distribuciones
dispersas puede ser sustancial.

\subsection{Intuicion geometrica}

La desigualdad de Jensen se puede ilustrar graficamente. La funcion $\ln x$
es cóncava: la recta que une dos puntos cualesquiera de la curva queda por
debajo de la curva. El valor esperado de $\ln u_i$ (que es un promedio ponderado
de $\ln u_i$) queda por debajo del logaritmo del valor esperado $\ln\E[u_i]$.
En el caso extremo donde $u_i$ es constante ($\Var(u_i) = 0$), la igualdad
se cumple: $\E[\ln u_i] = \ln \E[u_i]$.

\begin{clave}
  Por la desigualdad de Jensen aplicada a la funcion cóncava $\ln(\cdot)$:
  \[
    \E[\ln u_i] < \ln \E[u_i] = \ln 1 = 0
  \]
  La media del logaritmo del error es \textbf{estrictamente negativa} cuando
  $\E[u_i] = 1$. Cuantificadamente: $\E[\ln u_i] = -\sigma^2/2$ bajo log-normalidad.
\end{clave}

\section{Inciso (c): como hacer que $\E(\ln u_i) = 0$}

\subsection{Planteamiento del problema}

Suponga que $\E[\ln u_i] = w \neq 0$ (donde $w < 0$ por el resultado del
inciso (b) si $\E[u_i] = 1$). El modelo logaritmico~(\ref{eq:log}) es:
\begin{equation}
  \ln Y_i = \ln\beta_1 + \beta_2 \ln X_i + \underbrace{\ln u_i}_{=\,w + v_i}
  \label{eq:problema}
\end{equation}

\subsection{Reparametrizacion mediante absorcion en el intercepto}

Se define un nuevo termino de error centrado:
\begin{equation}
  v_i = \ln u_i - w = \ln u_i - \E[\ln u_i]
  \label{eq:vi_def}
\end{equation}

Por construccion:
\begin{equation}
  \E[v_i] = \E[\ln u_i] - w = w - w = 0
  \label{eq:vi_media_cero}
\end{equation}

Sustituyendo $\ln u_i = w + v_i$ en~(\ref{eq:problema}):
\begin{align}
  \ln Y_i
  &= \ln\beta_1 + \beta_2 \ln X_i + w + v_i
  \notag\\
  &= \underbrace{(\ln\beta_1 + w)}_{\displaystyle\alpha_1^*}
   + \beta_2 \ln X_i + v_i
  \label{eq:modelo_corregido}
\end{align}

El modelo~(\ref{eq:modelo_corregido}) cumple ahora $\E[v_i] = 0$ y puede
estimarse por MCO. El nuevo intercepto $\alpha_1^* = \ln\beta_1 + w$ absorbe
el componente no nulo de la media del error.

\subsection{Recuperacion del parametro original $\beta_1$}

Una vez estimado $\hat\alpha_1^*$, el parametro $\beta_1$ del modelo en
niveles~(\ref{eq:nivel}) se recupera como:
\begin{equation}
  \hat\beta_1 = e^{\hat\alpha_1^* - w} = e^{\hat\alpha_1^*} \cdot e^{-w}
  \label{eq:beta1_recuperado}
\end{equation}

Sin embargo, para aplicar~(\ref{eq:beta1_recuperado}) se necesita conocer $w$.
Si $w$ es desconocido (como suele ocurrir en la practica), se estima mediante:
\begin{equation}
  \hat w = \frac{1}{n}\sum_{i=1}^n \ln\hat u_i
  \label{eq:w_estimado}
\end{equation}
donde $\ln\hat u_i = \ln Y_i - \hat\alpha_1^* - \hat\beta_2 \ln X_i$ son
los residuos del modelo logaritmico. Nótese que si MCO se aplica correctamente
con un intercepto libre, las condiciones normales garantizan que la suma de
residuos es cero, por lo que $\hat w \equiv 0$ en el modelo estimado.

\subsection{Interpretacion alternativa}

La reparametrizacion muestra que el intercepto estimado $\hat\alpha_1^*$
no corresponde directamente a $\ln\hat\beta_1$, sino a $\ln\hat\beta_1 + w$.
Si se desea recuperar $\beta_1$ a partir del intercepto estimado, se necesita
aplicar la correccion de sesgo:
\begin{equation}
  \hat\beta_1^{\text{corregido}}
    = e^{\hat\alpha_1^*} \cdot e^{-\hat w}
    = e^{\hat\alpha_1^*}
      \cdot \exp\!\left(-\frac{\hat\sigma^2}{2}\right)
  \label{eq:beta1_corregido}
\end{equation}
donde $\hat\sigma^2$ es el estimador de la varianza de $\ln u_i$ bajo log-normalidad.
Esta correccion es conocida como \textbf{correccion de sesgo de Goldberger} (1968).

\begin{clave}
  Si $\E[\ln u_i] = w \neq 0$, se redefine $v_i = \ln u_i - w$, que satisface
  $\E[v_i] = 0$. El modelo transformado es:
  \[
    \ln Y_i = (\ln\beta_1 + w) + \beta_2 \ln X_i + v_i
  \]
  El nuevo intercepto $\alpha_1^* = \ln\beta_1 + w$ absorbe la media no nula
  del error. La pendiente $\beta_2$ no se afecta. Para recuperar $\beta_1$
  se debe aplicar la correccion $\hat\beta_1 = e^{\hat\alpha_1^*} \cdot e^{-\hat\sigma^2/2}$
  bajo log-normalidad.
\end{clave}

\section{Sintesis: propiedades de los modelos logaritmicos}

\begin{table}[ht]
\centering
\caption{Resumen de las propiedades del termino de error en modelos log-log}
\label{tab:resumen}
\small
\begin{tabular}{p{4cm} p{5cm} p{4.5cm}}
\toprule
\textbf{Supuesto} & \textbf{Consecuencia} & \textbf{Implicacion practica} \\
\midrule
$\ln u_i \sim \mathcal{N}(0, \sigma^2)$
  & $u_i \sim \text{Log-Normal}(0, \sigma^2)$; $\E[u_i] = e^{\sigma^2/2} > 1$
  & MCO sobre modelo log es valido si el intercepto es libre \\[6pt]
$\E[u_i] = 1$ (modelo multiplicativo)
  & $\E[\ln u_i] = -\sigma^2/2 < 0$ por desigualdad de Jensen
  & El intercepto del modelo log esta sesgado hacia abajo \\[6pt]
$\E[\ln u_i] = w \neq 0$
  & Redefine $v_i = \ln u_i - w$; $\E[v_i] = 0$
  & El intercepto absorbe $w$; la pendiente $\beta_2$ no se sesga \\[6pt]
Recuperacion de $\beta_1$
  & $\hat\beta_1 = e^{\hat\alpha_1^* - \hat w}$
  & Requiere correccion de sesgo de Goldberger (1968) \\
\bottomrule
\end{tabular}
\end{table}

\appendix

\section{Demostracion de la desigualdad de Jensen para $\ln(\cdot)$}

Sea $g(x) = \ln x$ para $x > 0$. Expansion de Taylor de segundo orden alrededor
de $\mu = \E[u_i]$:
\begin{equation}
  \ln u_i \approx \ln\mu + \frac{u_i - \mu}{\mu} - \frac{(u_i - \mu)^2}{2\mu^2}
  \label{eq:taylor_ln}
\end{equation}
Tomando esperanza:
\begin{equation}
  \E[\ln u_i]
    \approx \ln\mu + 0 - \frac{\Var(u_i)}{2\mu^2}
    = \ln\E[u_i] - \frac{\Var(u_i)}{2[\E[u_i]]^2}
  \label{eq:jensen_taylor}
\end{equation}
Como $\Var(u_i) > 0$ y $[\E[u_i]]^2 > 0$, el segundo termino es negativo, confirmando
$\E[\ln u_i] < \ln\E[u_i]$. Bajo log-normalidad con $\mu = e^{\mu_0 + \sigma^2/2}$
(media de la log-normal) y $\E[u_i] = 1$ (es decir, $\mu_0 = -\sigma^2/2$):
\begin{equation}
  \E[\ln u_i] = \mu_0 = -\frac{\sigma^2}{2}
  \label{eq:resultado_exacto}
\end{equation}
Este resultado es exacto (no una aproximacion) bajo log-normalidad.

\section{Correccion de sesgo de Goldberger (1968)}

En el modelo multiplicativo $Y_i = \beta_1 X_i^{\beta_2} u_i$ con
$u_i \sim \text{Log-Normal}$ tal que $\E[u_i] = 1$ (condicion de media uno),
el predictor MCO del logaritmo $\widehat{\ln Y_i} = \hat\alpha_1^* + \hat\beta_2 \ln X_i$
subestima sistematicamente $\E[\ln Y_i | X_i]$. El predictor correcto de $Y_i$
(no de $\ln Y_i$) es:

\begin{equation}
  \hat Y_i = e^{\hat\alpha_1^* + \hat\beta_2\ln X_i} \cdot e^{\hat\sigma^2/2}
  \label{eq:predictor_Goldberger}
\end{equation}

El factor $e^{\hat\sigma^2/2}$ es la correccion de Goldberger: sin el, el
predictor $e^{\hat\alpha_1^* + \hat\beta_2\ln X_i}$ subestima $\E[Y_i | X_i]$
en el factor $e^{-\hat\sigma^2/2}$.

\vspace{14pt}
\noindent\rule{\linewidth}{0.4pt}

\noindent{\small\textit{Nota.} El problema del sesgo al retransformar predicciones
logaritmicas a escala original es comun en econometria aplicada. Ademas de la
correccion de Goldberger para el caso log-normal, existen correcciones mas
generales (Duan, 1983) que no requieren suponer una distribucion especifica para
$u_i$. En la practica, muchos aplicadores ignoran este sesgo, lo que puede producir
predicciones significativamente subestimadas cuando $\hat\sigma^2$ es grande.}

\end{document}
